% SEGMENT OLP-0716-S01
% Part: second-order-logic
% Chapter: syntax-and-semantics
% Section: language-of-sol

\documentclass[../../../include/open-logic-section]{subfiles}

\begin{document}

\olfileid{sol}{syn}{lan}

\olsection{இரண்டாம்தரத் தருக்கத்தின் மொழி}

முதல்தரத் தருக்கத்தைப் போலவே, இரண்டாம்தரத் தருக்கத்தின் கோவைகள்
\emph{!!{variable}s}, \emph{!!{constant}s}, \emph{!!{predicate}s},
சில நேரங்களில் \emph{!!{function}s} ஆகியவற்றைக் கொண்ட அடிப்படைச்
சொற்கோவையிலிருந்து கட்டியெழுப்பப்படுகின்றன. இவற்றுடன் தருக்க
இணைப்பிகள், அளவையடைகள், அடைப்புக்குறிகள் மற்றும் காற்புள்ளிகள் போன்ற
நிறுத்தற்குறிகளைச் சேர்த்து, \emph{உறுப்புகளும்}
\emph{!!{formula}s உம்} அமைக்கப்படுகின்றன. பொருள்களுக்கான மாறிகளோடு,
உறவுகளுக்கும் சார்புகளுக்குமான மாறிகளையும் இரண்டாம்தரத் தருக்கம்
கொண்டிருப்பதும், அவற்றின்மீது அளவையிடலை அனுமதிப்பதுமே வேறுபாடு.
எனவே இரண்டாம்தரத் தருக்கத்தின் தருக்கக் குறிகள் முதல்தரத்
தருக்கத்தின் குறிகளுடன் பின்வருவனவற்றையும் கொண்டவை:

% SEGMENT OLP-0716-S02
\begin{enumerate}
\item ஒவ்வோர் இட எண்ணிக்கை~$n$ க்கும் இரண்டாம்தர உறவு
  !!{variable}s இன் ஓர் !!{denumerable}s கணம்:
  $\Obj V_0^n$, $\Obj V_1^n$, $\Obj V_2^n$, \dots
\item இரண்டாம்தரச் சார்பு !!{variable}s இன் ஓர் !!{denumerable}s கணம்:
  $\Obj u_0^n$, $\Obj u_1^n$, $\Obj u_2^n$, \dots
\end{enumerate}

% SEGMENT OLP-0716-S03
முதல்தரத் தருக்கத்தில் !!{identity}~$\eq$ பொதுவாகச் சேர்க்கப்படுகிறது.
முதல்தரத் தருக்கத்தில், ஒரு மொழி~$\Lang{L}$ இன் தருக்கமல்லாத குறிகள்
சுவையான எதையும் வெளிப்படுத்த இன்றியமையாதவை. தருக்கமல்லாத குறிகள்
எதையும் பயன்படுத்தாத !!{sentence}s நிச்சயமாக உள்ளன; ஆனால்~$\eq$
மட்டும் கொண்டு சுவையான எதையும் கூறுவது கடினம். இரண்டாம்தரத்
தருக்கத்தில், உறவு மற்றும் சார்பு மாறிகள் வரம்பின்றிக் கிடைப்பதால்,
தருக்கமல்லாத குறிகளின் தனிப்பட்ட தொகுதி இல்லாமலேயே ஒரு முதல்தர
மொழியில் சொல்லக்கூடிய எதையும் நம்மால் சொல்ல முடியும்.

\end{document}
